% ==============================================================================
% reports/frontmatter/abstract_en.tex
% Abstract (English)
% ==============================================================================

\begin{otherlanguage}{english}
\chapter*{Abstract}
\markboth{Abstract}{Abstract}
\addcontentsline{toc}{chapter}{Abstract}

In information security, applying machine learning (ML) to Network Intrusion Detection Systems (NIDS) has garnered significant attention. Empirical studies using random data partitioning can report near-perfect F1-scores ($> 99\%$), but these scores may not reflect performance under evolving traffic patterns or a different evaluation protocol.

This course project report presents an empirical evaluation of brute force attack detection (\FTP\ and \SSH) on the standard CICIDS2017 Tuesday dataset. By implementing a time-based split combined with conservative boundary purging ($q=60\text{s}$) and embargo ($g=120\text{s}$) techniques, the evaluation examines performance on an SSH subtype absent from the benchmark model's fit data: models are fitted on morning \FTP\ attacks and evaluated on afternoon \SSH\ attacks. Importantly, the validation split contains \DataValidationSSH\ SSH samples and participates in model and hyperparameter selection; hence, this setup evaluates a subtype absent from the benchmark's model-fit data, rather than a fully closed end-to-end zero-shot pipeline.

Empirical results across three tree-based algorithms (Decision Tree, Random Forest, and XGBoost) demonstrate severe performance degradation on the unobserved attack subtype: in the primary benchmark (\texttt{time/with\_port}), the representative model (Random Forest) achieves an F1-score of \RFTestFOne\ with an SSH detection rate of \RFTestSSHRecall\ (capturing only \RFTestTP\ out of \DataTestSSH\ SSH flows), while Decision Tree and XGBoost yield an F1-score of \DTTestFOne. Model explainability via SHAP and decision tree inspection shows strong model reliance on service ports (\texttt{Destination Port} exhibits the highest mean absolute SHAP value, \SHAPRatioTopOneTwo\ times greater than the second feature). When service ports are removed via port ablation, the SSH detection rate increases slightly to $\RFTestSSHRecall\ \text{--} \RFTestWithoutPortSSHRecall$ (with the highest time-based F1-score reaching \RFTestWithoutPortFOne\ for Random Forest \texttt{without\_port}), indicating that while the destination port serves as a primary shortcut, packet-level statistical differences between FTP and SSH also contribute to the observed distribution shift.

Under stratified random splitting, models attain high F1-scores of $0.9823 \text{--} 0.9997$. These results are compared across two protocols that also differ in subtype composition and test-set distribution; within-campaign dependence is a plausible contributor to optimistic estimates, but direct label or feature leakage has not been established. The XGBoost case study shows that scikit-learn's \texttt{refit=True} behaves as designed: the independent tuning run fits on Train + Validation, including \DataValidationSSH\ validation SSH flows. Because this pipeline also differs from the benchmark in features, preprocessing, and model configuration, the score difference cannot be attributed to refit alone.

\vspace{0.8cm}
\noindent\textbf{Keywords:} Network Intrusion Detection Systems (NIDS), CICIDS2017, Brute Force Attacks, Temporal Split, Optimistic Evaluation Bias, Cross-Subtype Generalization, Model Explainability (SHAP).
\end{otherlanguage}
